\chapter{From Experiment to Deployed, Monitored System}
\label{ch:projects-b}

\section{Phase 3 --- Experimentation Framework}

\begin{keytakeaway}
Random hyperparameter tuning is not experimentation. Real experimentation is: form a hypothesis
$\rightarrow$ design an experiment to test it $\rightarrow$ run it $\rightarrow$ analyze results $\rightarrow$
update your understanding $\rightarrow$ repeat.
\end{keytakeaway}

\subsection{Step 1: Establish a Baseline}

Try simple baselines before anything complex: statistical (mean/median prediction, last value for time series,
moving average, simple heuristic rules), business-rule baselines (the current manual process, expert-defined
thresholds, an existing system), and simple ML baselines (linear/logistic regression, a default-parameter
decision tree). If a complex neural network only beats linear regression by 2\%, question whether the complexity
is worth it; if your model doesn't beat ``predict last week's value,'' you have a problem. Document baseline
performance on \emph{all} your metrics, not just accuracy.

\subsection{Step 2: Hypothesis-Driven Experimentation}

Stop random searching. Use an explicit template for every experiment: \textbf{hypothesis} (``adding weather
features will improve demand forecasting because customer behaviour changes with weather''); \textbf{prediction}
(``RMSE should decrease by at least 10\% when weather features are included''); \textbf{design} (control: model
without the feature; treatment: model with it; everything else held constant); \textbf{results} (did it improve,
by how much, on which segments, any unexpected effects); \textbf{conclusion} (was the hypothesis supported, what
did you learn, what's next).

Experiment categories: feature experiments (adding/removing groups, encodings, interactions), model experiments
(family comparison, complexity, ensembling), data experiments (more data vs.\ better features, split strategy,
class balancing), hyperparameter experiments (impact, sensitivity, optimal ranges). Each experiment should answer
one specific question.

\subsection{Step 3: Experiment Tracking}

Track, for every experiment: code version (git commit hash), data version (DVC hash/snapshot), full configuration
(not just hyperparameters --- split parameters, feature choices, preprocessing, seeds), results (all metrics,
per-segment breakdowns, confusion matrix, feature importance), computational cost (training time, memory, GPU/CPU
count, inference latency), qualitative observations, and artifacts (model file, test predictions, visualizations,
error analysis). Tools: MLflow, Weights \& Biases, Neptune, or a spreadsheet for simple tracking. The goal: six
months later, you can reproduce any experiment exactly.

\subsection{Step 4: Model Selection Criteria}

Choose on performance (does it clear the business threshold, with confidence intervals considered --- not just
``higher is better''), interpretability (regulated industries, debugging, stakeholder trust all depend on it),
complexity (lines of code, dependencies, specialist knowledge, debugging difficulty), computational cost (training
cadence, inference latency vs.\ SLA, infrastructure cost), robustness (sensitivity to shift, missing features,
edge cases), and development speed (iteration and debugging speed).

\begin{keytakeaway}
The right model is the simplest one that meets requirements --- not the most accurate, not the most sophisticated.
\end{keytakeaway}

\subsection{Step 5: Error Analysis}

Understanding failures is more valuable than celebrating successes. Identify failure patterns (where does the
model perform worst; are errors random or systematic; do they cluster in segments); quantify their business cost
and priority; investigate root causes (data quality, missing features, model capacity, wrong inductive bias);
generate improvement hypotheses from the patterns found.

\begin{example}
\textbf{Worked example.} ``Model predicts demand poorly for new products (cold start). Possible solutions:
category-level features instead of product-specific ones; content-based features from product attributes; a
conservative default prediction for new items; a separate model for new vs.\ established products.'' This shows
you don't just tune hyperparameters --- you think about problems.
\end{example}

\section{Phase 4 --- Model Development \& Validation}

\subsection{Step 1: Train/Validation/Test Split Strategy}

\textbf{Time-based splits for time series} --- never use random splits when temporal dependencies exist. Typical
structure: training on the oldest 60\%, validation on the next 20\% (for tuning), test on the most recent 20\%
(final evaluation only) --- because future data will look like the test set, not the training set.
\textbf{Cross-validation for non-temporal data} --- k-fold for robust estimates, but consider stratified splits
for imbalanced classes, group splits where data has clusters (same user, same location), and leave-one-out for
very small datasets. The key rule: your validation strategy must simulate production reality --- validate on new
users if production will see new users; validate across seasons if production has seasonal patterns.

\subsection{Step 2: Evaluation Metrics Selection}

\begin{longtable}{p{0.3\textwidth} p{0.64\textwidth}}
\toprule
\textbf{Problem type} & \textbf{Appropriate metrics} \\
\midrule
\endhead
Classification, balanced, equal costs & Accuracy \\
Classification, imbalanced & Precision/Recall/F1, AUC-ROC, Average Precision \\
Classification, asymmetric costs & Cost-sensitive metrics, expected-value calculation \\
Multi-class & Macro vs.\ micro F1, confusion-matrix analysis \\
Regression, outliers matter & RMSE \\
Regression, outliers less important & MAE \\
Regression, percentage errors matter & MAPE \\
Regression, direction matters & Directional accuracy \\
Ranking & NDCG, MAP, MRR \\
\bottomrule
\end{longtable}

Never rely on just one metric. Example suite: primary = business metric (revenue impact); secondary = ML metric
(F1); monitoring = precision (limit false alarms) and recall (catch important cases); fairness = demographic
parity (equal error rates across groups).

\subsection{Step 3: Cross-Validation Strategy, Rigorously}

Check variance across folds (wildly different scores signal an unstable model or heterogeneous data); identify
and investigate the worst-performing fold; validate across relevant dimensions (temporal, spatial, demographic);
confirm statistical significance of any improvement over baseline with bootstrap confidence intervals or a formal
test. The goal is to convince yourself --- and recruiters --- that the model actually works, not that it got
lucky on one split.

\subsection{Step 4: Robustness Testing}

Test how the model handles: \textbf{missing features} (drop each one and measure the impact; identify critical
dependencies; build graceful degradation); \textbf{distributional shift} (different time periods, different
segments, simulated shifts via resampling); \textbf{adversarial inputs} (deliberately crafted edge cases, boundary
conditions, vulnerable features); \textbf{data-quality degradation} (injected noise, missingness, corrupted
values). Production always reveals weaknesses you didn't test for --- the more you test beforehand, the fewer
surprises later.

\subsection{Step 5: Model Interpretation}

Interpretation matters for debugging (understanding why a model failed), trust (stakeholders won't use
predictions they don't understand), compliance (regulated industries require explainability), and improvement
(you can't improve what you don't understand). \textbf{Global} methods: feature importance, partial-dependence
plots, model summary statistics. \textbf{Local} methods: SHAP values, LIME, counterfactual explanations.
\textbf{Model-specific}: coefficient analysis for linear models; split importance and tree visualization for tree
models; attention weights and activation maps for neural networks. Document findings concretely --- e.g.\ ``model
primarily relies on recency features (70\% of importance); day-of-week adds a 20\% bump on weekends; weather
contributes only 5\%, possibly low signal or poor encoding.''

\section{Phase 5 --- Deployment Architecture}

\begin{keytakeaway}
A model working in a Jupyter notebook means nothing. Production deployment requires reliability (graceful failure
handling), scalability (works under load), maintainability (others can debug and update it), and observability
(you know what's happening).
\end{keytakeaway}

\subsection{Step 1: Deployment Pattern Selection}

\begin{longtable}{p{0.16\textwidth} p{0.2\textwidth} p{0.1\textwidth} p{0.1\textwidth} p{0.32\textwidth}}
\toprule
\textbf{Pattern} & \textbf{When} & \textbf{Latency} & \textbf{Cost} & \textbf{Typical use cases} \\
\midrule
\endhead
Batch & Daily forecasts, weekly recs, monthly reports & Minutes--hours & Low & Demand forecasting, segmentation,
risk scoring, content recs \\
Real-time API & User-facing, instant response needed & $<$100\,ms & Medium & Fraud detection, personalization,
search ranking, moderation \\
Streaming & Event-driven, continuous processing & Seconds & High & Real-time analytics, anomaly detection, live
recs \\
Edge & Offline capability, privacy requirements & $<$10\,ms & Dev-heavy, low runtime & Mobile apps, IoT,
privacy-sensitive apps \\
\bottomrule
\end{longtable}

Most student projects should start with batch or a simple API --- don't overcomplicate without a real requirement.

\subsection{Step 2: Serving Architecture Design}

Components to consider: \textbf{model serving layer} (input validation, feature transformation, inference, output
formatting); \textbf{feature store}, only if the same features feed multiple models, feature computation is
expensive, or features must stay consistent across training/serving --- don't build one unless you actually need
it; \textbf{caching layer} for repeated, deterministic, stable queries; \textbf{load balancer}, only needed at
scale; \textbf{model registry} for version management, rollback, and A/B testing. Start simple (\texttt{User
$\rightarrow$ API $\rightarrow$ validation $\rightarrow$ feature engineering $\rightarrow$ prediction
$\rightarrow$ response}) and add complexity only when justified.

\subsection{Step 3: API Design Best Practices}

Define a clear input schema (required/optional fields, type/range/format validation) and output schema
(prediction, confidence, explanation where possible, metadata such as model version and timestamp). Handle errors
explicitly (400 for bad input, 500 for server errors, 429 for rate limiting, with helpful messages). Document
endpoints, request/response examples, authentication, rate limits, and error codes. Version the API so multiple
model versions can be supported during a transition. Add health-check and metrics endpoints. Apply basic security:
authentication, input sanitization, rate limiting. Even student projects benefit from getting these details right
--- it signals you understand production requirements.

\subsection{Step 4: Containerization Strategy}

Containerize for reproducibility (works on any machine with Docker), isolation (no dependency conflicts with the
host), portability (deploy anywhere), and version control (images are versioned artifacts). Include application
code, model artifacts, pinned dependencies, environment configuration, and a health-check script. Do \emph{not}
include training data, development tools, or secrets (use environment variables instead).

\subsection{Step 5: Deployment Checklist}

\begin{multicols}{2}
\begin{itemize}[nosep]
  \item Input validation implemented
  \item Error handling in place
  \item Logging configured
  \item Monitoring set up
  \item Documentation complete
  \item Rollback plan defined
  \item Testing done (unit, integration, load)
  \item Security reviewed
  \item Cost estimated
  \item Maintenance plan exists
\end{itemize}
\end{multicols}

Deployment is not ``push to production and hope.'' It's systematic preparation for things going wrong.

\section{Phase 6 --- Monitoring \& Maintenance Strategy}

\begin{keytakeaway}
If you can't measure it, you can't improve it. If you don't monitor it, it will break silently.
\end{keytakeaway}

\subsection{Step 1: What to Monitor}

\begin{description}
  \item[Business metrics] Revenue impact, cost savings, time saved, user satisfaction, conversion rates.
  \item[Model performance metrics] Prediction accuracy, precision/recall, RMSE/MAE, AUC/F1.
  \item[Data quality metrics] Missing-value rates, out-of-range values, feature distributions, schema violations.
  \item[System health metrics] Request latency, error rates, throughput, resource utilization.
\end{description}

\subsection{Step 2: Data Drift Detection}

\begin{longtable}{p{0.22\textwidth} p{0.38\textwidth} p{0.32\textwidth}}
\toprule
\textbf{Drift type} & \textbf{What changes} & \textbf{Detection} \\
\midrule
\endhead
Covariate shift (feature drift) & Input distribution shifts, relationship stays the same (user behaviour, market
conditions, seasonality) & Compare feature distributions, training vs.\ production \\
Concept drift & The relationship between features and target changes (market dynamics, preference shifts) &
Monitor prediction error rates over time \\
Label shift & Target distribution changes, relationship stays the same (class balance shifts, rare events become
common) & Track prediction-distribution changes \\
\bottomrule
\end{longtable}

Monitoring approach: statistical tests (KS test, chi-square, population stability index); thresholds (alert if a
feature mean shifts $>2\sigma$, prediction distribution shifts significantly, or error rate rises $>10\%$);
visualization (distribution plots over time, control charts, drift dashboards). On detection: investigate the
cause, retrain with recent data, adjust feature engineering, update the baseline. Don't just detect drift --- have
a response plan.

\subsection{Step 3: Performance Monitoring}

Ground truth arrives via user feedback, business outcomes, or expert review --- often with a delay, so track
prediction quality on a lag, use proxy metrics in the meantime, and watch which segments degrade first.
Dashboards: real-time (requests/sec, latency, error rate, confidence distribution), daily (accuracy where
available, feature-importance changes, drift scores, business-metric impact), weekly (performance trends, data
quality trends, reliability, cost). Alerts: \textbf{critical} (error rate $>5\%$, latency $>500$\,ms for $>10\%$
of requests, downtime, pipeline failure --- immediate action); \textbf{warning} (accuracy drop $>5\%$, drift
detected, unusual patterns, resource spikes --- investigate soon); \textbf{informational} (retrain recommended,
new data available, report ready --- review when convenient). The goal: know when something breaks before users
complain.

\subsection{Step 4: Retraining Strategy}

\textbf{Scheduled retraining} (fixed cadence) suits data that updates regularly and degrades predictably.
\textbf{Triggered retraining} (on degradation or drift) suits unpredictable data changes or expensive training. A
\textbf{hybrid} combines a regular schedule with performance-triggered emergency retrains. Pipeline: collect new
data $\rightarrow$ quality checks $\rightarrow$ feature engineering (same transforms) $\rightarrow$ train a
candidate $\rightarrow$ evaluate against the current production model $\rightarrow$ staged rollout (small \% of
traffic, monitor, ramp to 100\% if better; investigate, alert, and keep the current model if worse) $\rightarrow$
rollback capability retained throughout. The retraining pipeline should be fully automated, with humans involved
only for approval and monitoring.

\subsection{Step 5: Incident Response Plan}

\begin{longtable}{p{0.28\textwidth} p{0.64\textwidth}}
\toprule
\textbf{Incident type} & \textbf{Response} \\
\midrule
\endhead
Data pipeline failure (source unavailable, quality-check failure, processing error) & Use cached/fallback data;
alert the data team; document the issue \\
Model performance degradation (sudden accuracy drop, drift, adversarial inputs) & Investigate the cause; trigger
an emergency retrain; consider rollback \\
System outage (API downtime, infrastructure failure, dependency breakage) & Activate the fallback system; alert
the infrastructure team; run a post-mortem after resolution \\
Security incident (unauthorized access, breach, malicious input) & Isolate the system; investigate scope;
implement fixes; notify stakeholders \\
\bottomrule
\end{longtable}

Post-incident, always perform a root-cause analysis, an impact assessment, define prevention measures, and
document the incident and its learnings. Even student projects benefit from thinking through failure modes --- it
signals maturity and production thinking.

